\documentclass[Thesis.tex]{subfiles}
\begin{document}

\chapter{Introduction}

\section{Discrete conformal maps and uniformization}

The notion of discrete conformal equivalence of triangle meshes was introduced in \cite{Springborn2008} and extended to the hyperbolic case in \cite{Bobenko2010}. Two combinatorially equivalent euclidean triangle meshes are conformally equivalent if there exist scale factors associated to vertices such that corresponding edge lengths are equal up to multiplication with adjacent scale factors.
The problem of finding a discrete conformal map from a triangulated surface to the plane can be formulated in terms of a variational principle. 

\begin{figure}
\centering		
\includegraphics[width=0.25\linewidth]{introduction/genus0_image.png}\\
\includegraphics[width=0.25\linewidth]{introduction/genus0_domain.png}
\caption{Uniformization of compact Riemann surfaces}
\label{fig:intro_uniformization}
\end{figure}


\begin{theorem}{Uniformization}

\end{theorem}


This work is divided into three parts. Part~\ref{part:uniformization} contains joint work with Alexander Bobenko and Boris Springborn on Discrete uniformization of Riemann surfaces \cite{BobSechSpr}. 

\section{Applications}

\section{Software implementation}


\subfilebibliography
\end{document}

%%% Local Variables:
%%% TeX-master: "Thesis.tex"
%%% End: